%!TEX root = ../slides.tex
\section{HTTP}


\begin{frame}{HyperText Transfer Protocol}
  \begin{description}
		\item[HyperText] Text with links.
		\item[Transfer] Communication of data.
		\item[Protocol] Set of rules for communication.\\[1cm]
    \item[HTTP/1.0] was the first version.
    \item[HTTP/1.1] is the current widely used version.
    \item[HTTP/2] use is slowly growing.
  \end{description}
\end{frame}


\begin{frame}{Why study HTTP?}
  \begin{description}
		\item[HTTP] is the main protocol used by web browsers.
    \item[Netflix] uses HTTP to stream video (using DASH).
    \item[Instagram] is basically just a HTTP API.
    \item[Facebook] is a web application using HTTP.
    \item[GMail] uses HTTP.
    \item[Twitter] uses HTTP.
    \item[Medium] is really clever about its use of HTTP.
    \item[Sandvine] suggest most internet traffic happens over HTTP.  
  \end{description}
  \citeurl{sandvine.com/downloads/general/global-internet-phenomena/2015}
\end{frame}


\begin{frame}{Why is HTTP so widely used?}
  \begin{description}
		\item[HTTP] is often used instead of protocols that are more suited to the application.
    \item[Browsers] are one of the main reasons for this. Modern operating systems come with one or more browsers installed by default.
    \item[Web servers] and browsers mainly talk over HTTP.
    \item[Libraries] exist for most programming languages to make HTTP requests.
    \item[HTTP] is relatively straght-forward.
    \item[Firewalls] usually do not block HTTP by default.  
  \end{description}
\end{frame}


\begin{frame}{What is HTTP?}
  \begin{description}
		\item[You] will know HTTP from typing \mintinline{http}{http://} in your browser's location bar.
    \item[RFC 2616] details HTTP/1.1.
    \item[HTTP] is a standard way of transmitting data at the application level.
    \item[Originally] it was for transmitting text.
    \item[Text] is strings of ones and zeroes chopped into chunks, where each chuck represents a letter or character.
    \item[Nothing] stops us from sending non-text data via HTTP. 
  \end{description}
\end{frame}


\begin{frame}{How does HTTP work?}
  \begin{description}
		\item[Clients] performs requests. Firefox is an example of a client.
    \item[Servers] respond to requests. Apache is an example of a server.
    \item[Requests] are text documents sent by clients to servers.
    \item[Responses] are text documents sent by server to clients.
    \item[URLs] specify the server's location and the resource on the server.
    \item[Headers] are text metadata added to the start of requests and responses.
    \item[Bodies] are the main content of responses, and sometimes requests.
    \item[HTML] typically forms the body of a response.  
  \end{description}
\end{frame}

\begin{frame}{Request--Response}
  \tikzstyle{rect} = [rectangle,fill=gmitblue,text width=4.5em,text centered,minimum height=4em,rounded corners,text=white]
  \tikzstyle{line} = [draw,->,very thick]
  \tikzstyle{oval} = [ellipse,fill=gmitred,text width=5em,text centered,text=white]
  \begin{adjustbox}{max width={0.9\textwidth},center} 
    \begin{tikzpicture}[node distance = 4cm]
      \node [oval] (client1) {Client (Firefox)};
      \node [rect, right of=client1] (request) {Request};
      \node [rect, right of=client1] (request) {Request};
      \node [oval, right of=request] (server1) {Server (Apache)};
      \node [rect, below=1cm of server1] (generate) {Generate response};
      \node [oval, below=1cm of generate] (server2) {Server (Apache)};
      \node [rect, left of=server2] (response) {Response};
      \node [oval, left of=response] (client2) {Client (Firefox)};
      \path [line] (client1) -- (request);
      \path [line] (request) -- (server1);
      \path [line,dashed] (server1) -- (generate);
      \path [line,dashed] (generate) -- (server2);
      \path [line] (server2) -- (response);
      \path [line] (response) -- (client2);
    \end{tikzpicture}
  \end{adjustbox}
\end{frame}


\begin{frame}{HTTP is not like a phone call}
  \begin{description}
		\item[Suppose] I ring you and ask for your PPS Number.
    \item[If] you don't understand the question, you can say ``Sorry, can you please repeat that?''.
    \item[Then] I repeat my question, and you then give me the number before hanging up.\\[1cm]
    \item[HTTP] doesn't work like that.
    \item[Misunderstandings] result in the server responding with an error, and hanging up.
    \item[Status codes] indicate errors, amongst other things.
  \end{description}
\end{frame}


\begin{frame}{Uniform Resource Locator}
  \textcolor{BlueViolet}{http}://%
  \textcolor{RubineRed}{username}:%
  \textcolor{Mahogany}{password}@%
  \textcolor{MidnightBlue}{www}.%
  \textcolor{OliveGreen}{reddit.com}:%
  \textcolor{Dandelion}{80}%
  \textcolor{Plum}{/r/funny/}?%
  \textcolor{DarkOrchid}{limit=1}
  
  \begin{table}
    \begin{tabular}{r@{\hspace{0.5cm}}p{6cm}}
      \textcolor{BlueViolet}{http} & Protocol \\
      \textcolor{RubineRed}{username} & Username \\
      \textcolor{Mahogany}{password} & Password \\
      \textcolor{MidnightBlue}{www} & Subdomain \\
      \textcolor{OliveGreen}{reddit.com} & Domain \\
      \textcolor{Dandelion}{80} & Port \\
      \textcolor{Plum}{/r/funny/} & Path \\
      \textcolor{DarkOrchid}{limit=1} & Parameter
    \end{tabular}
  \end{table}
\end{frame}


\begin{frame}[fragile]{Request and Response Format}
  Requests and responses both have this format:
  \begin{itemize}
    \item Intial line.
    \item Zero or more header lines.
    \item A blank line.
    \item Optional message body (e.g. a HTML file)
  \end{itemize}
  \citeurl{www.jmarshall.com/easy/http}
\end{frame}


\begin{frame}[fragile]{Request (GET) Example}
  \begin{minted}{http}
GET /courses/all-courses HTTP/1.1
Host: gmit.ie
User-Agent: curl/7.50.1
Accept: */*

  \end{minted}
\end{frame}

\begin{frame}[fragile]{Response Example}
  \begin{minted}{http}
HTTP/1.1 200 OK
Date: Mon, 27 Jul 2009 12:28:53 GMT
Server: Apache/2.2.14 (Win32)
Last-Modified: Wed, 22 Jul 2009 19:15:56 GMT
Content-Length: 88
Content-Type: text/html
Connection: Closed

<html>
  <body>
    <h1>Hello, World!</h1>
  </body>
</html>
  \end{minted}
\end{frame}


\begin{frame}{Resources}
 \begin{columns}[onlytextwidth]
   \begin{column}{0.5\textwidth}
     \centering
      \begin{figure}
        \begin{tikzpicture}
          \draw (0,0) -- (0,1.2) -- (0.7,1.2) -- (0.7,0.8) -- (1,0.8) -- (1,0) -- cycle;
          \draw (0.7,1.2) -- (1,0.8);
        \end{tikzpicture}
        \caption*{File}
      \end{figure}
    \end{column}
    \begin{column}{0.5\textwidth}
      \centering
      \begin{figure}
      \begin{tikzpicture}
        \node (A) [cylinder, shape border rotate=90, draw,minimum height=1.3cm,minimum width=1cm] {};
      \end{tikzpicture}
      \caption*{Database}
      \end{figure}
    \end{column}
  \end{columns}
  
  \begin{quote}
    HTTP is used to transmit resources \ldots A resource is some chunk of information that can be identified by a URL \ldots The most common kind of resource is a file, but a resource may also be a dynamically-generated query result \ldots
  \end{quote}
  \citeurl{www.jmarshall.com/easy/http}
\end{frame}


\begin{frame}{HTTP Methods}
  \begin{description}
    \item[GET] Retrieve information from the server.
    \item[HEAD] Like get, but retrieve only the response header.
    \item[POST] Send data to the server.
    \item[PUT] Set the resource at the URL to the request data.
    \item[DELETE] Delete the resource at the URL.
    \item[CONNECT] Set up tunnel for other traffic to pass through HTTP.
    \item[OPTIONS] Find the allowable operations at the given URL.
    \item[TRACE] Echo the received request.
    \item[PATCH] Partial resource modification.
  \end{description}
\end{frame}


\begin{frame}{Status codes}
  \begin{description}
    \item[404] is probably the most famous status code. It means you requested a resource that doesn't exist on the server.
    \item[200] means everything is OK, and is the most common one in everyday browsing.
    \item[All] status codes are three digit numbers.
    \item[1xx] indicates an informational message only
    \item[2xx] indicates success of some kind
    \item[3xx] redirects the client to another URL
    \item[4xx] indicates an error on the client's part
    \item[5xx] indicates an error on the server's part
  \end{description}
  \citeurl{www.jmarshall.com/easy/http}
\end{frame}


\begin{frame}{Sending data to the server}
  \begin{description}
    \item[Sometimes] we want to tell the server something extra.
    \item[Resources] can be generated differently based on this extra data.
    \item[Google] can use this to perform searches on google.ie
    \item[Try] opening \mintinline{http}{google.ie} in your browser, and then open \mintinline{http}{google.ie/?q=gmit}.
    \item[The requested] resource here is the same but we don't get the same response.
    \item[Two] main ways to send extra data to the server: in the URL after a question mark, or in the body of the request.
  \end{description}
  \citeurl{www.jmarshall.com/easy/http}
\end{frame}


\begin{frame}{URL encoding}
  HTML form data is usually URL-encoded by changing:
  \begin{itemize}
    \item Unsafe characters to \% \emph{xx} where \emph{xx} is the ASCII value.
    \item All spaces to plusses.
    \item Names and values to: name1=value1\&name2=value2.
  \end{itemize}
  \hr
  \begin{description}
    \item[GET] --- in the URL after a question mark.
    \item[POST] --- in the body.
  \end{description}
  \citeurl{www.jmarshall.com/easy/http}
\end{frame}

\begin{frame}[fragile]{Request (GET) with parameters}
  \begin{minted}{http}
GET /r/ireland?limit=1 HTTP/1.1
Host: reddit.com
User-Agent: curl/7.50.1
Accept: */*

  \end{minted}
\end{frame}


\begin{frame}[fragile]{Request (POST) Example}
  \begin{minted}{http}
POST /path/script.cgi HTTP/1.0
From: frog@jmarshall.com
User-Agent: HTTPTool/1.0
Content-Type: application/x-www-form-urlencoded
Content-Length: 32

home=Cosby&favorite+flavor=flies
  \end{minted}
\end{frame}
